\documentclass[12pt]{article}

% Packages
\usepackage[utf8]{inputenc}
\usepackage[margin=1in]{geometry}
\usepackage{mathptmx} % Times New Roman font
\usepackage{graphicx}
\usepackage{setspace}
\usepackage{titlesec}
\usepackage{tocloft}
\usepackage{hyperref}
\usepackage{amssymb} % For checkmark

% Font size commands
\newcommand{\fourteenpt}{\fontsize{14}{16.8}\selectfont}
\newcommand{\thirteenpt}{\fontsize{13}{15.6}\selectfont}
\newcommand{\twelvept}{\fontsize{12}{14.4}\selectfont}

% Section formatting
\titleformat{\section}
  {\normalfont\fourteenpt\bfseries}{\thesection}{1em}{}
\titleformat{\subsection}
  {\normalfont\thirteenpt\bfseries}{\thesubsection}{1em}{}
\titleformat{\subsubsection}
  {\normalfont\twelvept\bfseries}{\thesubsubsection}{1em}{}

% Line spacing
\onehalfspacing

\begin{document}

% Title Page
\begin{titlepage}
\centering
\vspace*{2cm}

{\Large School of Science and Engineering}

\vspace{2cm}

{\Large \textbf{Containerized Workflow Orchestration for Polyglot Processes}}

\vspace{1cm}

{\Large Capstone Design Interim Report}

\vspace{0.5cm}


\vspace{2cm}

{\Large November 2025}

\vspace{2cm}

{\Large El Mehdi Laidoudi}

\vspace{1cm}

{\Large Supervised by:}

\vspace{0.5cm}

{\Large Dr. Naeem Nisar Sheikh}

\end{titlepage}

% Student Conduct and Copyright Page
\newpage
\section*{Student conduct and copyright}

The supervisor and the student (the Capstoner) agree that:

\begin{enumerate}
    \item Permission has been obtained for any third-party content (eg Data, illustrations, photographs, charts or maps).
    \item The results described in this report have not previously been published
\end{enumerate}

\vspace{2cm}

Student's name and signature: El Mehdi Laidoudi \hspace{2cm} \underline{\hspace{5cm}}

\vspace{2cm}

Supervisor's name and signature: \hspace{4cm} \underline{\hspace{5cm}}


\vspace{5cm}

\begin{center}
Capstone Report
\end{center}

\vspace{0.5cm}

\noindent
\textbf{Student Statement:}

I, El Mehdi Laidoudi, hereby declare that this report is an original account of the work I have done on my capstone project for Fall 2025. This document details my work on research and development of my project on building a containerized workflow platform.

\vspace{2cm}

\noindent
\makebox[\textwidth]{\hrulefill}

\begin{center}
Student's name:

El Mehdi Laidoudi
\end{center}

\vspace{3cm}

\begin{center}
Approved by the Supervisor
\end{center}

\vspace{2cm}

\noindent
\makebox[\textwidth]{\hrulefill}

\begin{center}
Dr. Naeem Nisar Sheikh
\end{center}

% Acknowledgements
\newpage
\section*{ACKNOWLEDGEMENTS}
\addcontentsline{toc}{section}{ACKNOWLEDGEMENTS}

\noindent
\parbox{\textwidth}{%
I would like to express my sincere gratitude to Dr. Naeem Nisar Sheikh for his invaluable guidance, support, and mentorship throughout this capstone project. His expertise and insights have been instrumental in shaping the direction of this work. I am also grateful to Al Akhawayn University for providing guidance. Special thanks to my family and friends for their unwavering support and encouragement throughout this journey. Finally, I acknowledge the open-source community whose tools and frameworks, including FastAPI, React Flow, and PostgreSQL, made this project possible.
}

% Table of Contents
\newpage
\tableofcontents

% List of Figures
\newpage
\listoffigures
\addcontentsline{toc}{section}{LIST OF FIGURES}

% List of Tables
\newpage
\listoftables
\addcontentsline{toc}{section}{LIST OF TABLES}

% Abstract
\newpage
\section*{ABSTRACT (ENGLISH)}
\addcontentsline{toc}{section}{ABSTRACT}

\noindent
\parbox{\textwidth}{%
This interim report presents the ongoing development of Orbitium, a containerized workflow orchestration platform being designed for polyglot scientific computing processes. The platform addresses the growing complexity of scientific workflows that require seamless integration of multiple programming languages and computational tools. Currently, we have implemented an intuitive visual interface for designing and managing computational workflows where each task can be implemented in different programming languages including Python, JavaScript, Go, and Ruby. The system architecture comprises a React-based frontend utilizing React Flow for visual workflow composition, a FastAPI backend for workflow management, and PostgreSQL for persistent storage. The workflow design phase has been completed, with nodes represented as directed acyclic graphs (DAGs). This interim report documents the progress to date, including system architecture, implemented features, and planned work for the execution orchestration phase which remains under development.
}

\vspace{1cm}

\noindent
\textbf{Keywords:} Workflow Design, Scientific Computing, Visual Programming, Polyglot Programming, Work in Progress

\vspace{2cm}

\section*{RESUME (FRENCH)}

\noindent
\parbox{\textwidth}{%
Ce rapport intermédiaire présente le développement en cours d'Orbitium, une plateforme d'orchestration de flux de travail conteneurisée en cours de conception pour les processus de calcul scientifique polyglotte. Actuellement, nous avons implémenté une interface visuelle intuitive pour la conception et la gestion de flux de travail computationnels où chaque tâche peut être implémentée dans différents langages de programmation. L'architecture du système comprend une interface frontend basée sur React, un backend FastAPI, et PostgreSQL pour le stockage persistant. La phase de conception des flux de travail a été achevée. Ce rapport intermédiaire documente les progrès réalisés à ce jour et les travaux planifiés pour la phase d'orchestration d'exécution qui reste en développement.
}

\vspace{1cm}

\noindent
\textbf{Mots clés:} Conception de Flux de Travail, Calcul Scientifique, Programmation Visuelle, Travaux en Cours

% Main Content
\newpage
\section{INTRODUCTION}

\hspace{\parindent}Scientific computing has become increasingly complex, requiring researchers to integrate multiple programming languages, frameworks, and computational tools to accomplish their objectives \cite{cohen2020scientific}. Modern scientific workflows often involve data ingestion, preprocessing, analysis, modeling, and visualization stages, each potentially requiring different technological stacks and execution environments \cite{deelman2015future}. This heterogeneity creates significant challenges for researchers who must manage dependencies, ensure reproducibility, and orchestrate execution across diverse computational resources.

\hspace{\parindent}Traditional workflow management systems often impose constraints on programming language choice or require extensive configuration and technical expertise \cite{amstutz2016common}. Researchers face the burden of containerizing individual components, managing inter-process communication, handling data flow between tasks, and ensuring proper execution sequencing. These technical obstacles divert valuable time and cognitive resources away from core scientific inquiry.

\hspace{\parindent}Orbitium addresses these challenges by providing a unified platform for visual workflow composition and orchestration that embraces polyglot programming paradigms. The system allows scientists to design computational workflows using an intuitive drag-and-drop interface where each node represents a discrete computational task that can be implemented in the programming language most appropriate for that specific operation. Table \ref{tab:languages} lists the currently supported programming languages.

\begin{table}[h]
\centering
\caption{Supported Programming Languages}
\label{tab:languages}
\begin{tabular}{|l|l|p{6cm}|}
\hline
\textbf{Language} & \textbf{Version} & \textbf{Primary Use Cases} \\
\hline
Python & 3.13 & Data analysis, machine learning, scientific computing \\
JavaScript & ES2023 & Data processing, web APIs, JSON manipulation \\
Go & 1.21+ & High-performance computing, concurrent operations \\
Ruby & 3.3 & Text processing, scripting, automation \\
\hline
\end{tabular}
\end{table}

\hspace{\parindent}By abstracting the complexity of containerization, dependency management, and execution orchestration, Orbitium enables researchers to focus on the scientific logic of their workflows rather than infrastructure concerns.

\hspace{\parindent}This capstone project presents the design, implementation, and validation of Orbitium, demonstrating how modern web technologies, containerization, and workflow orchestration principles can be synthesized to create an accessible platform for scientific computing. The system leverages React Flow for visual workflow design, FastAPI for robust backend services, PostgreSQL for data persistence, and Docker for containerized execution. The architecture supports directed acyclic graph (DAG) topology for workflow representation, ensuring proper task dependencies and enabling efficient parallel execution where possible.

\hspace{\parindent}The remainder of this interim report is structured as follows: Section 2 articulates the problem statement, Section 3 details current project specifications, Sections 4-6 analyze sociotechnical considerations, Section 7 reviews related literature, Section 8 describes the methodology and design approach taken so far, Section 9 presents preliminary results from the implemented components, Section 10 reflects on learning strategies employed, Section 11 outlines remaining work, and Section 12 provides interim conclusions.

\newpage
\section{PROBLEM STATEMENT}

\hspace{\parindent}Scientific research increasingly demands computational workflows that integrate multiple programming languages and tools, yet existing workflow management systems present significant barriers to adoption and effective use. Researchers face three interconnected problems that impede their computational work.

\hspace{\parindent}First, \textbf{technological fragmentation} forces scientists to work with disparate tools that don't interoperate seamlessly. A typical data science workflow might require Python for statistical analysis, R for specialized visualization, Julia for high-performance numerical computing, and shell scripts for data preprocessing. Existing systems either lock researchers into a single programming ecosystem or require extensive manual configuration to bridge between languages \cite{perkel2019workflow}. This fragmentation increases cognitive load and reduces productivity as researchers must context-switch between different development environments, dependency managers, and execution paradigms.

\hspace{\parindent}Second, \textbf{reproducibility challenges} plague scientific computing. Research must be reproducible to be credible, yet computational workflows often fail to reproduce due to dependency conflicts, environment inconsistencies, and insufficient documentation of execution contexts \cite{stodden2016enhancing}. While containerization technologies like Docker provide solutions, they require specialized knowledge that many researchers lack. The gap between container expertise and domain expertise creates a reproducibility crisis where valuable scientific findings cannot be reliably validated or extended.

\hspace{\parindent}Third, \textbf{complexity of orchestration} presents a formidable technical barrier. Coordinating task execution order, managing data flow between computational steps, handling failures gracefully, and parallelizing independent operations require sophisticated understanding of workflow management principles \cite{liu2015survey}. Researchers must manually encode task dependencies, write custom scripts for data passing, implement error handling logic, and optimize execution plans. This complexity diverts intellectual effort from scientific problem-solving to software engineering concerns.

\hspace{\parindent}These problems compound when researchers attempt to collaborate or share their computational methods. Workflows developed by one researcher often fail to execute in another's environment. The lack of standardized, accessible platforms for polyglot workflow orchestration limits scientific productivity, impedes reproducibility, and creates barriers to open science practices. 

\hspace{\parindent}Orbitium addresses these problems by providing an integrated platform that abstracts infrastructure complexity while preserving flexibility in language choice, ensures reproducibility through containerization, and manages orchestration automatically through visual workflow composition.

\newpage
\section{PROJECT SPECIFICATIONS}

\hspace{\parindent}Orbitium is designed as a web-based workflow orchestration platform with the following technical and functional specifications. Table \ref{tab:tech_stack} summarizes the complete technology stack.

\begin{table}[h]
\centering
\caption{Complete Technology Stack}
\label{tab:tech_stack}
\begin{tabular}{|l|l|l|}
\hline
\textbf{Layer} & \textbf{Technology} & \textbf{Version} \\
\hline
\multicolumn{3}{|c|}{\textbf{Frontend}} \\
\hline
Framework & React & 19.1.1 \\
Type System & TypeScript & Latest \\
Build Tool & Vite & 4.x \\
Workflow Visualization & React Flow (@xyflow/react) & 12.8.5 \\
State Management & Zustand & Latest \\
Routing & TanStack Router & 1.131.48 \\
Data Fetching & TanStack Query & 5.89.0 \\
Code Editor & Monaco Editor & 4.7.0 \\
UI Components & Radix UI & Latest \\
Graph Layout & Dagre & 1.1.5 \\
\hline
\multicolumn{3}{|c|}{\textbf{Backend}} \\
\hline
Framework & FastAPI & 0.118+ \\
Language & Python & 3.13 \\
ORM & SQLModel & 0.0.27 \\
Validation & Pydantic & via FastAPI \\
Package Manager & UV & Latest \\
\hline
\multicolumn{3}{|c|}{\textbf{Database}} \\
\hline
Database & PostgreSQL & 17-alpine \\
\hline
\multicolumn{3}{|c|}{\textbf{Infrastructure}} \\
\hline
Containerization & Docker & Latest \\
Orchestration & Docker Compose & Latest \\
\hline
\end{tabular}
\end{table}

\textbf{Core Functionality:}
\begin{itemize}
    \item Visual workflow design interface enabling drag-and-drop composition of computational pipelines
    \item Support for polyglot programming with initial implementation supporting Python, JavaScript, Go, and Ruby
    \item Directed Acyclic Graph (DAG) representation of workflow topology ensuring proper task dependencies
    \item Real-time workflow validation preventing cyclic dependencies and ensuring execution feasibility
    \item Integrated code editor with syntax highlighting and language-specific features powered by Monaco Editor
    \item Workflow persistence enabling save, load, and modification of computational pipelines
    \item Execution orchestration with proper task sequencing based on dependency resolution
\end{itemize}

\textbf{System Architecture:}
\begin{itemize}
    \item \textit{Frontend}: React 19.1 with TypeScript for type-safe development
    \item \textit{Workflow Visualization}: React Flow (XyFlow) version 12.8 for node-based workflow composition
    \item \textit{State Management}: Zustand for efficient client-side state management
    \item \textit{Routing}: TanStack Router for type-safe navigation with file-based routing
    \item \textit{Backend API}: FastAPI 0.118+ providing RESTful endpoints for workflow CRUD operations
    \item \textit{Database}: PostgreSQL 17 with SQLModel for type-safe ORM operations
    \item \textit{Containerization}: Docker for isolated execution environments
    \item \textit{Development Tooling}: Vite for fast development builds, UV for Python dependency management
\end{itemize}

\textbf{Data Models:}
The system implements a Flow model with the following schema:
\begin{itemize}
    \item Unique identifier (UUID) for each workflow
    \item Name and description metadata
    \item JSON-serialized workflow data containing nodes and edges arrays
    \item Node data structure: id, type, position coordinates, and task-specific data (title, description, executable specification)
    \item Edge data structure: id, source node, target node, and edge type
    \item Ready flag indicating workflow validation status
\end{itemize}

\textbf{API Endpoints:}

\begin{table}[h]
\centering
\caption{REST API Endpoints}
\label{tab:api_endpoints}
\begin{tabular}{|l|l|p{6cm}|}
\hline
\textbf{Method} & \textbf{Endpoint} & \textbf{Description} \\
\hline
GET & /api/v1/flows & Retrieve all workflows \\
GET & /api/v1/flows/\{id\} & Retrieve specific workflow by ID \\
POST & /api/v1/flows/create-flow & Create new workflow \\
PUT & /api/v1/flows/\{id\}/edit & Update existing workflow \\
DELETE & /api/v1/flows/\{id\} & Delete workflow \\
GET & /api/v1/flows/\{id\}/execute & Execute workflow (future) \\
GET & /api/v1/health-check & Health check endpoint \\
\hline
\end{tabular}
\end{table}

\textbf{Performance Requirements:}
\begin{itemize}
    \item Sub-second response time for workflow load and save operations
    \item Support for workflows with up to 100 nodes
    \item Automatic layout computation using Dagre for optimal node positioning
    \item Real-time visual feedback during workflow composition
\end{itemize}

\textbf{Security Considerations:}
\begin{itemize}
\item CORS configuration for controlled frontend-backend communication
\item Environment-based configuration management separating development and production settings
\item Database connection pooling for efficient resource utilization
\end{itemize}

\newpage
\section{STEEPLE ANALYSIS}

\hspace{\parindent}The STEEPLE framework provides systematic analysis of external factors affecting Orbitium's development and deployment across Social, Technological, Economic, Environmental, Political, Legal, and Ethical dimensions.

\hspace{\parindent}\textbf{Social:} The scientific community increasingly recognizes the importance of computational reproducibility and open science practices. Orbitium aligns with this cultural shift by lowering barriers to workflow sharing and collaboration. However, adoption depends on overcoming resistance to workflow tools among researchers who prefer traditional scripting approaches. The platform must provide sufficient value to justify learning a new system while respecting diverse computational practices across scientific disciplines.

\hspace{\parindent}\textbf{Technological:} Rapid evolution in containerization technologies, serverless computing, and workflow orchestration systems creates both opportunities and challenges. Docker's maturity provides stable foundation for containerized execution, while emerging standards like the Common Workflow Language (CWL) \cite{amstutz2016common} inform design decisions. The JavaScript and Python ecosystems' vibrant development ensures continued availability of tools and libraries. However, technological debt from framework updates and API changes requires ongoing maintenance.

\hspace{\parindent}\textbf{Economic:} Cloud computing costs represent a significant consideration for workflow execution at scale. While development uses local resources, production deployment requires balancing computational power against budget constraints. The open-source nature of core dependencies (React, FastAPI, PostgreSQL) minimizes licensing costs. Economic viability depends on demonstrating sufficient productivity gains to justify infrastructure investments. Research institutions face budget pressures that make cost-effective tools particularly valuable.

\hspace{\parindent}\textbf{Environmental:} Computational workflows consume significant energy, particularly for data-intensive scientific computing. Orbitium's containerized approach enables resource optimization and efficient scheduling. However, the platform must consider carbon footprint of cloud infrastructure. Future development should incorporate energy-aware scheduling algorithms and support for heterogeneous computing resources including energy-efficient ARM processors. Scientific computing's environmental impact increasingly influences tool adoption decisions.

\hspace{\parindent}\textbf{Political:} Data sovereignty and computational resource access vary across international boundaries. Scientific collaborations often span multiple countries with different regulatory frameworks. Orbitium must accommodate diverse deployment scenarios from institutional on-premise servers to multi-region cloud deployments. Political considerations affect choices about data residency, access controls, and compliance requirements.

\hspace{\parindent}\textbf{Legal:} Intellectual property concerns arise when workflows incorporate proprietary algorithms or datasets. The platform must enable appropriate access controls and licensing specification. Open-source licensing (MIT or Apache 2.0) for Orbitium itself encourages adoption while protecting developer interests. Compliance with data protection regulations like GDPR affects system design, particularly regarding workflow metadata and execution logs.

\hspace{\parindent}\textbf{Ethical:} Reproducibility constitutes an ethical imperative in scientific computing. Orbitium addresses this through containerization and explicit workflow specification. However, the platform must guard against facilitating harmful applications. Ethical considerations include responsible computing resource usage, accessibility for researchers with disabilities, and preventing weaponization of scientific workflows. Transparency in algorithm implementation and data provenance supports ethical scientific practices.

\newpage
\section{ENGINEERING STANDARDS}

\hspace{\parindent}Orbitium's development adheres to established engineering standards ensuring quality, interoperability, and maintainability. Table \ref{tab:engineering_standards} summarizes the key standards applied.

\begin{table}[h]
\centering
\caption{Engineering Standards Applied to Orbitium}
\label{tab:engineering_standards}
\begin{tabular}{|l|p{9cm}|}
\hline
\textbf{Standard} & \textbf{Application in Orbitium} \\
\hline
ISO/IEC 25010 & Quality model for maintainability, usability, and reliability \\
REST API & Stateless HTTP-based API design with standard methods \\
OpenAPI 3.0 & Automatic API documentation generation via FastAPI \\
Semantic Versioning & Dependency version management and compatibility \\
W3C Web Standards & HTML5, CSS3, ECMAScript 2023, ARIA accessibility \\
OCI & Docker image and runtime specifications for containers \\
12-Factor App & Configuration through environment, stateless processes \\
SQL:2016 & PostgreSQL standard compliance for database operations \\
ACID Properties & Transaction atomicity, consistency, isolation, durability \\
OWASP Top 10 & Security best practices for web applications \\
CORS & Cross-Origin Resource Sharing for secure API access \\
TypeScript Strict & Enhanced type safety in frontend development \\
Python Type Hints & Type safety and static analysis in backend code \\
CWL Compatibility & Node-based structure enables future CWL export \\
\hline
\end{tabular}
\end{table}

\textbf{Software Engineering Standards:}
\begin{itemize}
    \item \textit{ISO/IEC 25010 Software Quality Model}: The system architecture addresses quality characteristics including functionality, reliability, usability, efficiency, maintainability, and portability. React's component-based architecture supports maintainability through modular design. FastAPI's automatic API documentation enhances usability.
    
    \item \textit{RESTful API Design Principles}: The backend follows REST architectural constraints with stateless client-server communication, uniform resource identification through URIs, and standard HTTP methods (GET, POST, PUT, DELETE) for CRUD operations. JSON serves as the standardized data interchange format.
    
    \item \textit{Semantic Versioning (SemVer)}: Dependencies follow semantic versioning enabling predictable updates and backward compatibility management. The project structure supports independent versioning of frontend and backend components.
\end{itemize}

\textbf{Web Standards:}
\begin{itemize}
    \item \textit{W3C Web Standards}: The frontend implements HTML5, CSS3, and ECMAScript 2023 standards ensuring cross-browser compatibility. ARIA attributes support accessibility requirements for users with disabilities.
    
    \item \textit{OpenAPI Specification}: FastAPI automatically generates OpenAPI 3.0 compliant API documentation enabling standardized API description and client generation. This facilitates integration with external tools and services.
\end{itemize}

\textbf{Containerization Standards:}
\begin{itemize}
    \item \textit{OCI (Open Container Initiative)}: Docker images conform to OCI image and runtime specifications ensuring portability across container runtimes. Multi-stage builds optimize image sizes following best practices.
    
    \item \textit{Twelve-Factor App Methodology}: The application design follows twelve-factor principles including codebase management, dependency declaration, configuration through environment variables, and stateless processes.
\end{itemize}

\textbf{Database Standards:}
\begin{itemize}
    \item \textit{SQL Standards (SQL:2016)}: PostgreSQL implements SQL standards for relational data management. SQLModel provides type-safe query construction preventing SQL injection vulnerabilities.
    
    \item \textit{ACID Properties}: Database transactions maintain Atomicity, Consistency, Isolation, and Durability ensuring data integrity during concurrent workflow operations.
\end{itemize}

\textbf{Security Standards:}
\begin{itemize}
    \item \textit{OWASP Top 10}: The application addresses common web vulnerabilities including injection attacks through parameterized queries, broken authentication through secure session management, and sensitive data exposure through HTTPS enforcement.
    
    \item \textit{CORS (Cross-Origin Resource Sharing)}: Proper CORS configuration prevents unauthorized cross-origin requests while enabling legitimate frontend-backend communication.
\end{itemize}

\textbf{Code Quality Standards:}
\begin{itemize}
    \item \textit{TypeScript Strict Mode}: Frontend code uses TypeScript strict mode for enhanced type safety and early error detection.
    
    \item \textit{Python Type Hints}: Backend code extensively uses type hints enabling static analysis through mypy and improved IDE support.
    
    \item \textit{Linting Standards}: Ruff for Python and ESLint for TypeScript enforce consistent code style and identify potential errors.
\end{itemize}

\textbf{Workflow Standards:}
\begin{itemize}
    \item \textit{DAG Representation}: Directed Acyclic Graph structure aligns with established workflow representation standards used by Apache Airflow, Luigi, and Prefect. This facilitates potential future interoperability.
    
    \item \textit{Common Workflow Language (CWL)}: While not directly implemented, the node-based structure enables future CWL export capability supporting workflow portability across platforms \cite{amstutz2016common}.
\end{itemize}

\newpage
\section{LOGIC MODEL FRAMEWORK}

\hspace{\parindent}The Logic Model provides a systematic framework for understanding Orbitium's development process, resource allocation, activities, and expected outcomes.

\subsection{TARGET POPULATION}

\hspace{\parindent}Orbitium targets three primary user segments within the scientific computing community:

\hspace{\parindent}\textbf{Primary Users - Research Scientists and Data Scientists:} Computational researchers across domains including bioinformatics, climate science, physics, economics, and social sciences who need to orchestrate multi-stage analytical workflows. These users possess domain expertise and programming proficiency but lack specialized knowledge in DevOps and workflow orchestration systems. They value tools that accelerate research productivity without requiring extensive infrastructure management.

\hspace{\parindent}\textbf{Secondary Users - Research Software Engineers:} Technical specialists who support research teams by building and maintaining computational infrastructure. These users possess strong software engineering skills and seek platforms that balance ease of use for researchers with sufficient flexibility for complex requirements. They require robust APIs, extensibility mechanisms, and integration capabilities.

\hspace{\parindent}\textbf{Tertiary Users - Students and Educators:} Graduate students learning computational methods and instructors teaching scientific computing courses. This population benefits from accessible workflow tools that demonstrate computational thinking principles and best practices in reproducible research. Educational contexts prioritize intuitive interfaces and clear conceptual models.

\subsection{UNDERLYING ASSUMPTIONS}

\hspace{\parindent}The project development rests on several key assumptions:

\textbf{Technical Assumptions:}
\begin{itemize}
    \item Users possess basic programming skills in at least one supported language
    \item Institutional infrastructure supports Docker containerization
    \item Network connectivity enables web-based application access
    \item Modern web browsers (Chrome, Firefox, Safari, Edge) are available
    \item PostgreSQL database can be deployed in target environments
\end{itemize}

\textbf{User Behavior Assumptions:}
\begin{itemize}
    \item Scientists prefer visual workflow design over code-only approaches for certain tasks
    \item Reducing technical barriers increases adoption of workflow best practices
    \item Reproducibility concerns motivate investment in workflow tooling
    \item Researchers will document workflows adequately when tools make documentation easy
\end{itemize}

\textbf{Environmental Assumptions:}
\begin{itemize}
    \item Research institutions continue investing in computational infrastructure
    \item Open-source workflow tools remain viable against commercial alternatives
    \item Container technologies maintain community support and development
    \item Scientific community values interoperability and open standards
\end{itemize}

\subsection{RESSOURCES/CHALLENGES}

\textbf{Resources Utilized:}
\begin{itemize}
    \item \textit{Human Capital}: Developer time for architecture, implementation, testing, and documentation. Supervisor mentorship for guidance and feedback.
    
    \item \textit{Technical Infrastructure}: Development workstation, cloud hosting (potential deployment), GitHub for version control and collaboration.
    
    \item \textit{Software Resources}: Open-source frameworks (React, FastAPI, PostgreSQL), development tools (VS Code, Docker Desktop), documentation platforms.
    
    \item \textit{Knowledge Resources}: Academic literature on workflow systems, technical documentation for frameworks, online communities (Stack Overflow, GitHub Discussions).
    
    \item \textit{Financial Resources}: Minimal direct costs due to open-source stack. Potential cloud hosting costs for demonstration deployment.
\end{itemize}

\textbf{Key Challenges:}
\begin{itemize}
    \item \textit{Technical Complexity}: Integrating multiple technologies (React Flow, FastAPI, PostgreSQL, Docker) requires understanding diverse architectural patterns and ensuring smooth interoperation.
    
    \item \textit{Containerization Orchestration}: Implementing actual container execution for workflow nodes requires Docker API integration, resource management, and error handling beyond current prototype scope.
    
    \item \textit{State Management}: Synchronizing workflow state between frontend visualization, backend persistence, and execution engine presents consistency challenges.
    
    \item \textit{Performance Optimization}: Large workflows with many nodes require efficient rendering, layout computation, and data serialization strategies.
    
    \item \textit{User Experience}: Balancing simplicity for novice users with power features for advanced users requires careful interface design and progressive disclosure.
    
    \item \textit{Testing and Validation}: Comprehensive testing of visual workflow editor, API endpoints, and database operations demands significant effort.
    
    \item \textit{Documentation}: Creating user documentation, API references, and developer guides competes with implementation time.
\end{itemize}

\subsection{ACTIVITIES}

The project development encompasses several interconnected activities:

\textbf{Design Phase:}
\begin{itemize}
    \item System architecture planning defining frontend, backend, and database components
    \item User interface mockups and interaction flow design
    \item Database schema design for workflow representation
    \item API endpoint specification and documentation
    \item Technology stack evaluation and selection
\end{itemize}

\textbf{Implementation Phase:}
\begin{itemize}
    \item Frontend development: React components, state management with Zustand, React Flow integration
    \item Workflow visualization: Node types (start, task, end, initial), edge rendering, drag-and-drop functionality
    \item Code editor integration: Monaco Editor configuration, language support, syntax highlighting
    \item Backend development: FastAPI application structure, endpoint implementation, database integration
    \item Data persistence: SQLModel models, database migrations, CRUD operations
    \item Layout algorithms: Dagre integration for automatic workflow layout, position calculation
    \item API integration: Axios client, TanStack Query for state synchronization, error handling
\end{itemize}

\textbf{Testing Phase:}
\begin{itemize}
    \item Unit testing for utility functions and state management logic
    \item Integration testing for API endpoints and database operations
    \item End-to-end testing for workflow creation, editing, and saving
    \item User interface testing for visual consistency and interaction responsiveness
    \item Performance testing for workflow load times and rendering efficiency
\end{itemize}

\textbf{Documentation Phase:}
\begin{itemize}
    \item Code documentation with inline comments and docstrings
    \item API documentation through OpenAPI/Swagger interface
    \item User guide creation with screenshots and tutorials
    \item Technical report writing documenting design decisions and implementation
    \item README files for project setup and development instructions
\end{itemize}

\subsection{OUTPUTS OF THE PROJECT}

\textbf{Software Deliverables:}
\begin{itemize}
    \item Fully functional web application accessible through modern browsers
    \item Source code repository with version control history and commit documentation
    \item Docker Compose configuration for local development environment setup
    \item Database schema and migration scripts for PostgreSQL deployment
    \item REST API with automated OpenAPI documentation
    \item Containerized deployment artifacts (Dockerfiles, configuration files)
\end{itemize}

\textbf{Documentation Deliverables:}
\begin{itemize}
    \item Comprehensive technical report (this document) detailing system design and implementation
    \item User documentation explaining workflow creation and management
    \item Developer documentation for contributing and extending the platform
    \item API reference documentation with endpoint specifications and examples
    \item Architecture diagrams illustrating system components and interactions
\end{itemize}

\textbf{Demonstration Artifacts:}
\begin{itemize}
    \item Example workflows demonstrating platform capabilities
    \item Video demonstrations of workflow creation and execution processes
    \item Performance benchmarks and scalability analysis results
    \item Conference proposal for potential academic presentation
\end{itemize}

\subsection{OUTCOMES}

\textbf{Short-term Outcomes (Immediate):}
\begin{itemize}
    \item Successful capstone project completion demonstrating technical proficiency
    \item Working prototype validating core workflow orchestration concepts
    \item Portfolio artifact showcasing full-stack development capabilities
    \item Academic credit toward degree completion
    \item Enhanced understanding of workflow systems, containerization, and web technologies
\end{itemize}

\textbf{Medium-term Outcomes (6-12 months):}
\begin{itemize}
    \item Potential adoption by research groups for managing computational workflows
    \item Community contributions and feedback driving platform improvements
    \item Integration with additional programming languages and execution environments
    \item Publication opportunities at scientific computing or workflow conferences
    \item Foundation for continued open-source development and maintenance
\end{itemize}

\textbf{Long-term Outcomes (1-3 years):}
\begin{itemize}
    \item Established platform within scientific workflow ecosystem alongside Apache Airflow and Prefect
    \item Contribution to improving reproducibility practices in computational science
    \item Reduction in technical barriers enabling broader participation in computational research
    \item Potential commercial applications or institutional adoption at scale
    \item Influence on workflow management system design and best practices
\end{itemize}

\textbf{Impact Indicators:}
\begin{itemize}
    \item Number of workflows created and executed on the platform
    \item User satisfaction metrics and feedback quality
    \item GitHub repository stars, forks, and contributions
    \item Citations in academic literature or technical blogs
    \item Time savings reported by users compared to manual workflow management
    \item Reproducibility improvements in published computational research
\end{itemize}

\newpage
\section{LITERATURE REVIEW}

\hspace{\parindent}Scientific workflow management systems have evolved substantially over the past two decades, driven by increasing computational complexity and demands for reproducibility in scientific computing. This review examines relevant literature across workflow orchestration, containerization, visual programming, and polyglot computing paradigms.

\hspace{\parindent}\textbf{Workflow Management Systems:} Deelman et al. \cite{deelman2015future} provide comprehensive analysis of workflow management system evolution, identifying key challenges including scalability, fault tolerance, and resource optimization. Their work establishes taxonomy of workflow systems distinguishing between data-driven and control-driven approaches. Orbitium adopts control-driven paradigm with explicit task dependencies represented through DAG topology. Liu et al. \cite{liu2015survey} survey scientific workflow systems, highlighting the trade-off between generality and domain-specificity. While systems like Apache Airflow provide general-purpose orchestration, domain-specific tools offer optimized functionality for particular scientific domains. Orbitium positions itself as general-purpose with extensibility mechanisms for domain customization.

\hspace{\parindent}\textbf{Containerization and Reproducibility:} Boettiger \cite{boettiger2015introduction} introduces containerization as solution to reproducibility challenges in computational science, demonstrating how Docker enables consistent execution environments across diverse infrastructure. Stodden and Miguez \cite{stodden2016enhancing} extend this analysis, arguing containerization addresses but does not completely solve reproducibility challenges without workflow specification and provenance tracking. Orbitium integrates containerization with explicit workflow representation, addressing both environmental consistency and logical reproducibility. Gruening et al. \cite{gruening2018practical} examine practical implementation of containerized workflows in bioinformatics, providing design patterns that inform Orbitium's architecture.

\hspace{\parindent}\textbf{Common Workflow Language:} Amstutz et al. \cite{amstutz2016common} present Common Workflow Language (CWL) as standardization effort enabling workflow portability across execution platforms. CWL's declarative approach separates workflow logic from execution details, facilitating reproducibility and tool interoperability. While Orbitium currently uses proprietary workflow representation, the node-edge structure enables future CWL export capability. Crusoe et al. \cite{crusoe2022methods} evaluate CWL implementations, demonstrating benefits of standardized workflow specification for scientific collaboration.

\hspace{\parindent}\textbf{Visual Programming and Workflow Design:} Burnett and Baker \cite{burnett2016scaling} analyze scaling challenges in visual programming languages, identifying principles for managing complexity in node-based interfaces. Their cognitive dimensions framework informs Orbitium's interface design, particularly regarding viscosity (resistance to change) and visibility (ability to view workflow structure). Green and Petre \cite{green1996usability} examine usability analysis of visual programming environments, providing heuristics applied to React Flow integration. Their work emphasizes importance of secondary notation (visual cues beyond formal syntax) implemented through color coding and icon selection in Orbitium's node design.

\hspace{\parindent}\textbf{Polyglot Programming:} Meyerovich and Rabkin \cite{meyerovich2013empirical} conduct empirical study of programming language adoption, revealing pragmatic factors driving language choice including library availability, community support, and task appropriateness. Scientific computing particularly benefits from polyglot approaches as different languages excel at different computational tasks. Cohen et al. \cite{cohen2020scientific} argue for polyglot flexibility in scientific software, noting tension between language familiarity and computational efficiency. Orbitium addresses this by enabling per-node language selection within unified workflow.

\hspace{\parindent}\textbf{Graph Algorithms for Workflow:} Tarjan's \cite{tarjan1972depth} depth-first search algorithm provides foundation for topological sorting used in workflow dependency resolution. Orbitium implements topological sort for determining execution order from DAG representation. Kahn's \cite{kahn1962topological} alternative algorithm offers benefits for parallel execution planning, informing future optimization directions. Coffman and Graham \cite{coffman1972optimal} analyze optimal scheduling strategies for DAG-structured task systems, establishing theoretical bounds relevant to workflow execution optimization.

\hspace{\parindent}\textbf{Web-Based Scientific Applications:} Russell et al. \cite{russell2018workflow} examine design patterns for web-based scientific workflow tools, identifying requirements for responsiveness, real-time collaboration, and offline capability. Their findings inform Orbitium's architecture using React for reactive user interfaces and optimistic updates for perceived performance. Perkel \cite{perkel2019workflow} provides practitioner perspective on workflow tools in scientific research, highlighting importance of gentle learning curves and incremental adoption pathways.

\hspace{\parindent}\textbf{State Management in Complex Applications:} The Flux architecture \cite{flux2014facebook} introduces unidirectional data flow patterns addressing state management challenges in complex user interfaces. Zustand, employed in Orbitium, builds on these principles with simplified API reducing boilerplate. Redux documentation \cite{redux2015abramov} articulates principles of predictable state updates through immutability and pure functions, patterns adopted in Orbitium's workflow state management.

\hspace{\parindent}\textbf{API Design and Documentation:} Fielding's \cite{fielding2000architectural} REST architectural style establishes principles for scalable web APIs implemented in Orbitium's FastAPI backend. Richardson and Ruby \cite{richardson2008restful} provide practical guidance for RESTful API design, particularly regarding resource modeling and HTTP method semantics. OpenAPI Specification \cite{openapi2021specification} standardizes API documentation, automatically generated by FastAPI ensuring consistent interface contracts.

\hspace{\parindent}\textbf{Database Design for Scientific Data:} Date \cite{date2003database} provides foundational database theory informing Orbitium's relational schema design. The workflow representation balances normalized relational structure with JSON document storage for flexible node/edge data. Hellerstein and Stonebraker \cite{hellerstein2005readings} compile influential database papers establishing principles for transaction management and query optimization relevant to workflow metadata persistence.

\hspace{\parindent}\textbf{Gaps in Existing Literature:} While extensive literature addresses individual components of workflow systems, limited research examines integration of visual workflow design, polyglot programming support, and containerized execution in accessible web platforms. Existing systems typically require significant technical expertise or limit language flexibility. Orbitium contributes by demonstrating practical integration of these capabilities in user-centered design targeting scientific researchers rather than DevOps specialists.

\newpage
\section{METHODOLOGY \& CAPSTONE DESIGN}

\hspace{\parindent}This section details the technical methodology, architectural decisions, and implementation approach for Orbitium, organized into system design, data presentation, algorithm development, and software implementation.

\subsection{CAPSTONE DESIGN \& APPROACHES/ MODEL/ALGORITHM/APPLICATION DEVELOPMENT}

\hspace{\parindent}\textbf{System Architecture:} Orbitium implements a three-tier architecture separating presentation, application logic, and data persistence layers. The presentation tier uses React with TypeScript, providing type-safe component development and enhanced developer experience. React Flow (version 12.8.5) handles workflow visualization, offering node-based interface with built-in drag-and-drop, panning, and zooming capabilities. The component architecture follows atomic design principles with atoms (buttons, inputs), molecules (node components, edge components), organisms (workflow canvas, navigation bar), and templates (page layouts).

The application tier consists of FastAPI backend (version 0.118+) providing RESTful API endpoints. FastAPI's asynchronous request handling enables concurrent operations while automatic OpenAPI documentation generation ensures API discoverability. The backend implements CRUD operations for workflow management, validation logic for workflow integrity, and orchestration logic for execution coordination (future implementation). Dependency injection through FastAPI's Depends mechanism promotes testable, modular code.

The data persistence tier uses PostgreSQL 17 with SQLModel ORM providing type-safe database operations. SQLModel combines SQLAlchemy's maturity with Pydantic's validation, enabling shared models between API schemas and database tables. The database schema includes Flow table with UUID primary key, workflow metadata (name, description), and JSONB column storing workflow topology (nodes and edges arrays). JSONB enables flexible schema evolution while maintaining query capabilities.

\hspace{\parindent}\textbf{Workflow Representation Model:} Workflows are represented as directed acyclic graphs (DAGs) where nodes represent computational tasks and edges represent dependencies. This mathematical model ensures well-defined execution semantics: a workflow $W = (V, E)$ consists of vertex set $V$ representing tasks and edge set $E \subseteq V \times V$ representing dependencies. The acyclic constraint prevents infinite loops, ensuring workflows have finite execution time. Formal definition requires for any edge $(u, v) \in E$, there exists no path from $v$ back to $u$.

Each node $v \in V$ contains metadata: unique identifier, node type (start, task, initial, end), position coordinates $(x, y)$, and task data including title, description, and executable specification. The executable specification comprises programming language identifier and code string. Table \ref{tab:node_types} describes the different node types.

\begin{table}[h]
\centering
\caption{Workflow Node Types}
\label{tab:node_types}
\begin{tabular}{|l|l|p{6cm}|}
\hline
\textbf{Node Type} & \textbf{Dimensions} & \textbf{Purpose} \\
\hline
Start & 80 × 20 px & Triggers workflow execution, single source node \\
Task & 344 × 160 px & Computational operation with code and metadata \\
Initial & 208 × 120 px & Placeholder for workflow construction \\
End & 80 × 40 px & Marks workflow completion, single sink node \\
\hline
\end{tabular}
\end{table}

Node types implement different visualization and interaction patterns: start nodes trigger workflow execution, task nodes represent computational operations, initial nodes provide placeholders for workflow construction, and end nodes mark workflow completion.

Edges $e = (u, v) \in E$ connect source node $u$ to target node $v$, establishing that $v$ depends on $u$ and cannot execute until $u$ completes. Edge types control interaction behavior: addNodeEdge type enables inserting nodes between existing tasks, preserving workflow topology while supporting incremental construction.

\hspace{\parindent}\textbf{Layout Algorithm:} Orbitium employs Dagre (Data-driven Graph Rendering Engine) for automatic workflow layout. Dagre implements Sugiyama-style hierarchical graph drawing, organizing nodes in layers to minimize edge crossings and optimize visual clarity. The algorithm proceeds in four phases: 

(1) \textit{Layer assignment} using longest path algorithm assigns each node to layer based on maximum distance from source nodes. This ensures dependencies point consistently downward.

(2) \textit{Crossing reduction} applies heuristic methods to minimize edge crossings between adjacent layers, improving readability.

(3) \textit{Coordinate assignment} determines horizontal positions within layers, balancing edge length and spacing.

(4) \textit{Edge routing} computes smooth paths between nodes, implemented as smooth step curves in React Flow.

The implementation configures Dagre with rankdir='TB' (top-to-bottom layout), nodesep=50px (horizontal spacing), and ranksep=80px (vertical spacing). Node dimensions vary by type, encoded in getNodeDimensions utility function mapping types to width-height pairs: start nodes (80x20), end nodes (80x40), initial nodes (208x120), task nodes (344x160).

\hspace{\parindent}\textbf{State Management Architecture:} Zustand manages client-side application state with simplified API compared to Redux. The workflow store maintains nodes array, edges array, and nodeCounter for ID generation. Actions include loadFlow for initializing workflow from database, updateNodeMetaData for modifying task descriptions, setNodeBaseLanguage for changing programming language, updateNodeCode for code editing, and addNodeOnEdge for inserting tasks.

State updates follow immutability principles: modifications create new array copies rather than mutating existing state. This enables React's reconciliation algorithm to efficiently detect changes and re-render affected components. For example, updateNodeCode maps over nodes array, returning new objects for modified nodes while preserving references to unchanged nodes.

\hspace{\parindent}\textbf{Topological Sorting Algorithm:} Workflow execution requires topological ordering determining valid execution sequence. Algorithm \ref{alg:topo_sort} presents the implementation. The implementation uses depth-first search (DFS) approach:

\begin{verbatim}
function topologicalSort(nodes, edges):
    visited = set()
    order = []
    
    function dfs(node):
        if node in visited:
            return
        visited.add(node)
        for edge in edges where edge.source == node:
            dfs(edge.target)
        order.prepend(node)
    
    for node in nodes:
        if node not in visited:
            dfs(node)
    
    return order
\end{verbatim}

This algorithm has $O(|V| + |E|)$ time complexity, linear in workflow size. The prepend operation constructs reverse postorder traversal, ensuring each node appears after its dependencies. Cycle detection occurs naturally: if DFS encounters already-visited node in current path, a cycle exists and workflow is invalid.

\subsubsection{DATA PRESENTATION}

\hspace{\parindent}\textbf{Frontend Data Flow:} The application implements unidirectional data flow following Flux architecture principles. User interactions trigger actions that update Zustand store state. State changes propagate to subscribed React components, causing re-renders with updated data. React Flow receives nodes and edges arrays as props, rendering visual representation. User manipulations (dragging nodes, connecting edges) trigger callbacks that dispatch new actions, completing the cycle.

TanStack Query (React Query) manages server state with distinct lifecycle from UI state. The library provides hooks for data fetching (useQuery), mutations (useMutation), and cache management. Query keys implement hierarchical structure: flowKeys.all, flowKeys.lists(), flowKeys.detail(id) enable targeted cache invalidation. When workflow is saved, updateFlow mutation invalidates relevant queries, triggering refetch to synchronize with server state.

The separation of concerns ensures UI remains responsive during network operations. Optimistic updates immediately reflect user changes in interface before server confirmation, with automatic rollback on errors. Loading and error states receive explicit handling, providing feedback for asynchronous operations.

\hspace{\parindent}\textbf{API Data Contract:} FastAPI endpoints accept and return JSON payloads conforming to Pydantic models. Table \ref{tab:data_models} describes the data models.

\begin{table}[h]
\centering
\caption{Data Models and Schemas}
\label{tab:data_models}
\begin{tabular}{|l|l|p{6cm}|}
\hline
\textbf{Model} & \textbf{Type} & \textbf{Fields} \\
\hline
FlowBase & Schema & name (str), description (str), data (dict | null) \\
FlowRead & Response & FlowBase + id (UUID) \\
FlowUpdate & Request & name (optional), description (optional), data (optional) \\
Flow & Database & FlowBase + id (UUID, primary key), data (JSONB) \\
\hline
\multicolumn{3}{|c|}{\textbf{Workflow Data Structure}} \\
\hline
Node & Object & id, type, position \{x, y\}, data \\
NodeData & Object & title, nodeType, description, executable \\
Executable & Object & language (str), code (str) \\
Edge & Object & id, source (node id), target (node id), type \\
\hline
\end{tabular}
\end{table}

FlowBase schema defines creation payload: name (string), description (string), data (optional object containing nodes and edges). FlowRead schema adds id (UUID) for responses. FlowUpdate schema makes all fields optional, supporting partial updates.

The backend performs validation automatically through Pydantic, rejecting malformed requests with detailed error messages. Type coercion handles common discrepancies (string to integer conversion), while strict validation prevents invalid data from persisting. UUID type ensures globally unique workflow identifiers, eliminating collision risks in distributed deployments.

\hspace{\parindent}\textbf{Database Schema:} The Flow table implements following structure:

\begin{verbatim}
CREATE TABLE flow (
    id UUID PRIMARY KEY DEFAULT uuid_generate_v4(),
    name VARCHAR NOT NULL,
    description VARCHAR NOT NULL,
    data JSONB
);

CREATE INDEX idx_flow_name ON flow(name);
\end{verbatim}

JSONB type provides JSON storage with binary encoding, enabling efficient queries on workflow structure while maintaining flexibility. Indexes on name column accelerate workflow lookup by name. The schema supports future extensions: adding user\_id for multi-tenancy, created\_at/updated\_at for auditing, tags for categorization, without breaking existing code.

\hspace{\parindent}\textbf{Visualization Design:} Task nodes implement card-based design with gradient backgrounds, rounded corners, and subtle shadows creating depth perception. Each node displays: (1) language icon from stack-logos directory, (2) task title with bold font, (3) node type badge with colored background, (4) description text with line clamping, (5) handle components at top and bottom for edge connections.

Color coding provides visual distinction: teal-to-blue gradients for connection handles, gray gradients for node backgrounds, blue badges for node types. Monaco Editor integration provides professional code editing experience with syntax highlighting, auto-completion, and error detection for supported languages.

\subsubsection{EXECUTION ORCHESTRATION (PLANNED)}

\hspace{\parindent}The current implementation focuses on workflow design and persistence. Execution orchestration represents the next major development phase. The planned approach involves:

\hspace{\parindent}\textbf{Container Management:} Docker SDK for Python enables programmatic container lifecycle management. Each task node converts to container specification: base image determined by language (python:3.13-slim, node:20-alpine, golang:1.21-alpine, ruby:3.3-alpine), code mounted as volume, environment variables for configuration.

\hspace{\parindent}\textbf{Dependency Resolution:} Topological sort determines execution order. Tasks without dependencies execute immediately (potentially in parallel). Upon task completion, dependent tasks become eligible. Execution engine maintains ready queue of executable tasks and waiting set of tasks with unsatisfied dependencies.

\hspace{\parindent}\textbf{Data Passing:} Inter-task communication through volume mounts or message queues. Each task writes output to designated location, consumed by dependent tasks. JSON serialization enables language-agnostic data exchange. For large datasets, object storage (S3, MinIO) provides scalable alternative.

\subsubsection{SOFTWARE}

\hspace{\parindent}\textbf{Development Environment:} The project uses modern development tooling: VS Code as primary IDE with extensions for TypeScript, Python, Docker, and Git. UV (Python package manager) provides fast dependency resolution and virtual environment management. Vite serves as frontend build tool, offering instant hot module replacement and optimized production builds. Docker Desktop enables local containerization testing.

\hspace{\parindent}\textbf{Version Control:} Git manages source code with branching strategy: main branch for stable releases, development branch for integration, feature branches for new functionality. Commit messages follow conventional commits specification enabling automated changelog generation. GitHub hosts repository with issue tracking for bug reports and feature requests.

\hspace{\parindent}\textbf{Code Quality:} TypeScript strict mode catches type errors at compile time. ESLint enforces consistent code style and identifies anti-patterns. Prettier handles automatic code formatting, eliminating style debates. On backend, Ruff provides fast Python linting combining functionality of multiple tools (Flake8, isort, pyupgrade).

\hspace{\parindent}\textbf{Testing Strategy:} Unit tests validate individual functions and components using Vitest for frontend, pytest for backend. Integration tests verify API endpoint behavior with test database. End-to-end tests (planned) will use Playwright for browser automation, testing complete user workflows.

\hspace{\parindent}\textbf{Deployment Architecture:} Docker Compose orchestrates local development environment with services for database (PostgreSQL container), backend (FastAPI with hot reload), and frontend (Vite dev server). Environment variables manage configuration: database credentials, API endpoints, CORS settings. Production deployment (planned) uses container registry (Docker Hub, GitHub Container Registry) with orchestration platform (Docker Swarm, Kubernetes).

\hspace{\parindent}\textbf{Monitoring and Observability:} FastAPI provides automatic logging for requests and errors. Structured logging with JSON format enables log aggregation and analysis. Health check endpoint (/api/v1/health-check) enables uptime monitoring. Future enhancements include metrics collection (Prometheus), distributed tracing (OpenTelemetry), and error tracking (Sentry).

\newpage
\section{PRELIMINARY RESULTS AND CURRENT STATUS}

\hspace{\parindent}This section presents preliminary results from the implementation work completed to date. The workflow design and persistence components have been successfully implemented and tested, while the execution orchestration remains under development.

\hspace{\parindent}\textbf{Workflow Creation Performance (Preliminary):} Initial testing with workflows of varying complexity shows promising performance. Table \ref{tab:performance} summarizes preliminary metrics from the implemented features.

\begin{table}[h]
\centering
\caption{Preliminary Performance Metrics (Completed Features Only)}
\label{tab:performance}
\begin{tabular}{|l|r|l|}
\hline
\textbf{Feature} & \textbf{Metric} & \textbf{Status} \\
\hline
\multicolumn{3}{|c|}{\textbf{Workflow Design Interface}} \\
\hline
Small workflows (5-10 nodes) & $<$50 ms & \checkmark Tested \\
Medium workflows (20-30 nodes) & 100-150 ms & \checkmark Tested \\
Large workflows (50+ nodes) & TBD & \textit{Pending} \\
\hline
\multicolumn{3}{|c|}{\textbf{API Operations (Implemented)}} \\
\hline
GET /flows & 15-25 ms & \checkmark Tested \\
POST /create-flow & 20-30 ms & \checkmark Tested \\
PUT /\{id\}/edit & 25-35 ms & \checkmark Tested \\
GET /\{id\}/execute & N/A & \textit{Not Implemented} \\
\hline
\multicolumn{3}{|c|}{\textbf{Execution Engine}} \\
\hline
Container orchestration & N/A & \textit{Future Work} \\
Task scheduling & N/A & \textit{Future Work} \\
Parallel execution & N/A & \textit{Future Work} \\
\hline
\end{tabular}
\end{table}

Small workflows (5-10 nodes) load instantaneously with sub-50ms rendering time. Medium workflows (20-30 nodes) maintain responsive interaction with 100-150ms load time. Large workflows (50+ nodes) experience 200-300ms initial layout computation but remain interactive after loading. The Dagre layout algorithm exhibits $O(n \log n)$ complexity for $n$ nodes, confirming theoretical expectations.

Performance profiling using React DevTools identifies layout computation as primary bottleneck for large workflows. Optimization through memoization of layout results and incremental updates for node position changes maintains 60fps during user interactions. The React Flow library efficiently handles 100+ nodes through virtual rendering and viewport culling.

\hspace{\parindent}\textbf{Database Operation Latency:} API endpoint response times measured under local development conditions: GET /flows retrieves workflow list in 15-25ms, POST /create-flow creates new workflow in 20-30ms, PUT /{id}/edit updates workflow in 25-35ms, GET /{id} fetches specific workflow in 10-20ms. These measurements include network overhead, database query execution, and JSON serialization. PostgreSQL query execution remains consistently sub-10ms for workflow CRUD operations with current dataset scale.

JSONB storage demonstrates efficient serialization and deserialization. A workflow with 30 nodes serializes to approximately 15KB JSON payload, transmitted without perceptible delay over local network. Compression could reduce payload size by 60-70\% for production deployments over internet connections.

\hspace{\parindent}\textbf{Workflow Validation:} Cycle detection successfully prevents invalid workflow configurations. Testing with deliberately constructed cyclic workflows confirms detection algorithm identifies cycles and prevents persistence. The validation executes in $O(|V| + |E|)$ time, linear in workflow size, completing under 1ms for workflows up to 100 nodes.

Edge cases receive proper handling: workflows with disconnected subgraphs remain valid (multiple independent pipelines), workflows with single node represent valid minimal workflow, empty workflows (no computational nodes) are rejected as invalid. The system correctly handles concurrent editing scenarios through optimistic concurrency control, with last-write-wins semantics for conflicting updates.

\hspace{\parindent}\textbf{User Interface Responsiveness:} Interaction latency measurements show: node drag operations render at 60fps with no perceptible lag, edge creation provides immediate visual feedback, code editor updates reflect changes with <50ms latency after typing cessation, workflow save operation provides feedback within 100ms of user action.

The Monaco Editor integration performs well for code snippets up to 1000 lines, covering typical workflow task complexity. Syntax highlighting applies correctly for Python, JavaScript, Go, and Ruby. Auto-completion suggestions appear within 200ms of trigger character, comparable to dedicated IDEs.

\hspace{\parindent}\textbf{Cross-Browser Compatibility:} Testing across Chrome 120, Firefox 121, Safari 17, and Edge 120 confirms consistent rendering and functionality. React Flow's SVG rendering ensures visual fidelity across browsers. Minor styling discrepancies in rounded corners and gradients do not impact usability. Mobile responsiveness achieves basic functionality though touchscreen workflow editing requires interface optimization for practical use.

\hspace{\parindent}\textbf{Example Workflow Validation:} A representative scientific workflow implements temperature data analysis pipeline with four stages: (1) Load Temperature Data node (Python) reads CSV file, (2) Clean \& Validate Data node (Python) handles missing values and outliers, (3) Statistical Analysis node (Python) computes summary statistics, (4) Generate Visualization node (Python) creates plots. This workflow demonstrates end-to-end capability from design to persistence.

The workflow topology correctly maintains task dependencies through edge connections. Visual layout positions nodes in logical top-to-bottom flow with appropriate spacing. Code editor enables implementation of analysis logic in each node with syntax highlighting and validation. Workflow saves successfully to database and reloads with full state preservation including code content and node positions.

\hspace{\parindent}\textbf{State Synchronization:} Testing state consistency between frontend store, backend API, and database confirms reliable synchronization. Workflow modifications in UI update Zustand store immediately (optimistic update), propagate to backend API via PUT request within 100ms, persist to PostgreSQL within additional 20ms. Query invalidation triggers refetch confirming server state matches client state. Network failure scenarios correctly roll back optimistic updates, preventing inconsistent state.

\hspace{\parindent}\textbf{Scalability Considerations:} Linear scaling characteristics for workflow operations suggest system remains viable for hundreds of workflows and workflows with dozens of nodes. Database queries benefit from indexing on workflow name, with query plans showing efficient index scans. Current architecture does not address concurrent execution of multiple workflows, deferring scheduling concerns to future execution engine implementation.

\hspace{\parindent}\textbf{Code Quality Metrics:} Static analysis tools report: TypeScript strict mode identifies zero type errors in frontend code, ESLint reports no violations with current rule configuration, Ruff identifies no issues in Python backend code. Test coverage (limited in current prototype): critical utility functions achieve 80\%+ coverage, API endpoints covered by integration tests, React components have basic rendering tests.

\hspace{\parindent}\textbf{Memory Consumption:} Browser memory profiling shows: initial application load consumes ~50MB heap space, loading workflow with 30 nodes adds ~5MB, Monaco Editor instances contribute ~10MB per node with editor open. Memory remains stable during extended editing sessions with no observable leaks. Backend process consumes ~100MB resident memory under development server, primarily SQLAlchemy connection pool and FastAPI application object.

\hspace{\parindent}\textbf{Limitations and Constraints:} Current implementation limitations include: (1) no actual container execution implemented, workflow execution remains conceptual, (2) no user authentication/authorization, all workflows globally accessible, (3) no real-time collaboration features, concurrent editing may conflict, (4) limited error recovery mechanisms for database failures, (5) no workflow versioning or history tracking.

These limitations represent scope boundaries for capstone project rather than fundamental architectural constraints. The system design anticipates these features through appropriate abstraction layers and extension points.

\hspace{\parindent}\textbf{Current Progress Summary:} The preliminary results demonstrate successful implementation of the workflow design components: visual interface, polyglot node support, and persistence layer. The implemented features provide a solid foundation for the remaining work. The next phase will focus on execution engine integration, which represents the core innovation of containerized polyglot workflow execution.

\hspace{\parindent}\textbf{Challenges Encountered:} State synchronization between frontend and backend required careful architecture design. The graph layout algorithm needed optimization for larger workflows. These challenges have been addressed in the current implementation.

\hspace{\parindent}\textbf{Next Steps:} The execution orchestration phase will implement Docker container management, task scheduling, and inter-task communication. This represents approximately 40\% of the remaining development effort.

\newpage
\section{REMAINING WORK AND TIMELINE}

\hspace{\parindent}This section outlines the work remaining to complete the Orbitium platform and achieve the project objectives.

\subsection{EXECUTION ENGINE IMPLEMENTATION}

\hspace{\parindent}\textbf{Priority: High | Estimated Effort: 4-6 weeks}

The execution engine represents the most critical remaining component. Key tasks include:

\begin{itemize}
    \item Docker SDK integration for container lifecycle management
    \item Topological sort implementation for task scheduling
    \item Container specification generation from workflow nodes
    \item Inter-task data passing mechanisms (volume mounts or message queues)
    \item Error handling and recovery strategies
    \item Execution progress monitoring and logging
\end{itemize}

\subsection{TESTING AND VALIDATION}

\hspace{\parindent}\textbf{Priority: High | Estimated Effort: 2-3 weeks}

Comprehensive testing of the complete system:

\begin{itemize}
    \item Unit tests for execution engine components
    \item Integration tests for end-to-end workflow execution
    \item Performance testing with various workflow sizes
    \item Error scenario testing (container failures, network issues)
    \item Real scientific workflow validation cases
\end{itemize}

\subsection{DOCUMENTATION AND DEPLOYMENT}

\hspace{\parindent}\textbf{Priority: Medium | Estimated Effort: 1-2 weeks}

Finalizing project deliverables:

\begin{itemize}
    \item User documentation and tutorials
    \item API documentation completion
    \item Deployment guide for production environments
    \item Final report writing and presentation preparation
    \item Video demonstration of complete system
\end{itemize}

\subsection{OPTIONAL ENHANCEMENTS}

\hspace{\parindent}\textbf{Priority: Low | Time Permitting}

If time allows, these features would enhance the platform:

\begin{itemize}
    \item User authentication and authorization
    \item Workflow versioning and history
    \item Parallel execution optimization
    \item Real-time execution monitoring dashboard
    \item Workflow sharing and templates
\end{itemize}

\subsection{PROJECT TIMELINE}

\hspace{\parindent}\textbf{Weeks 1-6:} Execution engine implementation and core orchestration features

\hspace{\parindent}\textbf{Weeks 7-9:} Testing, debugging, and validation with scientific workflows

\hspace{\parindent}\textbf{Weeks 10-11:} Documentation, deployment preparation, final report

\hspace{\parindent}\textbf{Week 12:} Final presentation and project submission

\newpage
\section{LEARNING STRATEGIES}

\hspace{\parindent}This section reflects on learning methodologies employed throughout the capstone project and knowledge acquisition across technical and professional domains.

\subsection{TECHNICAL SKILL DEVELOPMENT}

\hspace{\parindent}\textbf{Full-Stack Development Competency:} The project necessitated simultaneous mastery of frontend and backend technologies. Frontend development with React and TypeScript required understanding component lifecycle, hooks API, state management patterns, and type system intricacies. Learning strategy combined official documentation study, tutorial completion, and iterative experimentation. TypeScript's strict typing initially increased cognitive load but ultimately improved code quality and refactoring confidence.

Backend development with FastAPI introduced asynchronous Python programming paradigms, dependency injection patterns, and ORM abstractions through SQLModel. The learning curve involved understanding async/await syntax, coroutine execution models, and database transaction management. Documentation exploration supplemented by Stack Overflow queries addressed specific implementation questions. Building API incrementally from simple endpoints to complex CRUD operations reinforced concepts through practical application.

\hspace{\parindent}\textbf{Graph Theory Application:} Implementing workflow orchestration required applying graph algorithms to practical problems. Topological sorting algorithms, previously understood theoretically, became concrete through implementation. Depth-first search traversal for cycle detection demonstrated algorithmic thinking applied to user-facing features. This reinforced connections between computer science fundamentals and software engineering practice.

Dagre layout library exposed hierarchical graph drawing algorithms and layer assignment strategies. Understanding algorithm selection trade-offs (speed vs. aesthetic quality) informed integration decisions. The experience highlighted importance of leveraging existing libraries while understanding underlying principles.

\hspace{\parindent}\textbf{Containerization and DevOps:} Docker integration required conceptual understanding of container technology, image layers, volume mounts, and networking. Initial confusion about container lifecycle and orchestration resolved through hands-on experimentation with docker-compose. The learning process involved iterative debugging of configuration issues, environment variable management, and database connection problems.

Understanding multi-stage builds, layer caching, and image optimization emerged from performance considerations. This practical DevOps knowledge complements theoretical software engineering coursework, preparing for industry expectations regarding deployment and infrastructure management.

\subsection{SOFTWARE ENGINEERING PRACTICES}

\hspace{\parindent}\textbf{Architecture and Design Patterns:} Designing system architecture required applying software design principles: separation of concerns, single responsibility, dependency inversion. Decisions about component boundaries, API design, and state management reflected engineering trade-offs between flexibility, complexity, and performance.

The iterative refinement process involved initial overengineering tendencies (premature abstraction) corrected through pragmatic simplification. Learning to balance ideal architectural patterns against project constraints (time, scope) developed professional judgment. Code reviews from supervisor provided valuable feedback improving design decisions.

\hspace{\parindent}\textbf{Version Control Proficiency:} Git usage evolved from basic commit operations to sophisticated workflows: feature branches for isolated development, meaningful commit messages documenting changes, strategic commits creating logical checkpoints. Understanding when to commit (atomic changes) vs. when to refactor before committing improved repository history quality.

Handling merge conflicts, using git stash for context switching, and writing descriptive commit messages became routine practices. The importance of version control for collaboration and experimentation risk mitigation became visceral through project experience.

\hspace{\parindent}\textbf{Testing and Quality Assurance:} Writing tests challenged initial focus on feature implementation. Learning test-driven development principles and experiencing benefits (refactoring confidence, documentation through examples) motivated investing in test coverage. Understanding test pyramid concepts (unit tests abundant, integration tests moderate, end-to-end tests minimal) informed testing strategy.

Static analysis tools (TypeScript compiler, ESLint, Ruff) provided automated quality checks, catching errors before runtime. Configuring these tools appropriately balanced strictness against productivity. The experience demonstrated value of tooling in maintaining code quality across growing codebase.

\subsection{RESEARCH AND PROBLEM-SOLVING}

\hspace{\parindent}\textbf{Documentation Navigation:} Effective use of technical documentation proved essential. Learning to navigate API references, identify relevant examples, and understand conceptual guides accelerated development. Distinguishing official documentation from community tutorials and evaluating source credibility became important skills.

Reading open-source code for libraries (React Flow, FastAPI) provided insights into professional code organization and implementation patterns. This "reading code to learn" approach complemented traditional learning methods.

\hspace{\parindent}\textbf{Debugging Methodology:} Systematic debugging practices developed through necessity: using browser developer tools for frontend debugging, logging statements for backend debugging, network tab inspection for API issues. Understanding stack traces, identifying error patterns, and searching error messages effectively reduced debugging time.

Rubber duck debugging and explaining problems to supervisor often revealed solutions through articulation process. This metacognitive strategy proved valuable for complex issues.

\hspace{\parindent}\textbf{Community Resources:} Stack Overflow, GitHub issues, and Discord communities provided problem-solving assistance. Learning to formulate good questions (minimal reproducible examples, clear problem statement) improved response quality. Contributing back through documentation improvements and issue reports fostered community engagement mindset.

\subsection{PROJECT MANAGEMENT}

\hspace{\parindent}\textbf{Scope Management:} Initial ambitious scope required pruning to achievable milestones. Learning to distinguish "must-have" from "nice-to-have" features prevented scope creep. Regular supervisor meetings provided accountability and feedback guiding priority decisions.

Creating task breakdowns in GitHub issues organized work into manageable chunks. Estimating effort for tasks improved with experience, revealing optimism bias in initial estimates.

\hspace{\parindent}\textbf{Time Management:} Balancing capstone work with coursework demanded effective time allocation. Blocking dedicated development time, setting weekly milestones, and tracking progress maintained momentum. Learning to switch between focused implementation and broad planning prevented both bikeshedding and direction loss.

\hspace{\parindent}\textbf{Documentation Discipline:} Maintaining documentation alongside development (rather than deferred to end) reduced cognitive load. Writing README content, API documentation, and code comments incrementally preserved context. This report writing process synthesized understanding and revealed knowledge gaps requiring additional learning.

\subsection{REFLECTIONS AND GROWTH}

\hspace{\parindent}The capstone project provided integration point for computer science education, connecting theoretical concepts (graph algorithms, database theory, software architecture) with practical implementation. The experience revealed software engineering as creative problem-solving requiring technical skills, design sensibility, and pragmatic judgment.

Challenges overcome (state management complexity, asynchronous programming confusion, deployment debugging) built confidence in tackling unfamiliar technical domains. The learning process itself—navigating documentation, seeking help appropriately, iterating on solutions—constitutes transferable meta-skill applicable to future technical challenges.

Future learning goals identified through project experience include: deeper understanding of distributed systems for scaling workflow execution, proficiency with Kubernetes for container orchestration, expertise in real-time web technologies for collaborative features, and formal methods for workflow verification.

\newpage
\section{INTERIM CONCLUSION}

\hspace{\parindent}This capstone project presents Orbitium, a containerized workflow orchestration platform designed to address challenges in scientific computing through visual workflow composition and polyglot programming support. The system successfully demonstrates integration of modern web technologies (React, FastAPI, PostgreSQL) with workflow management principles, providing accessible interface for researchers to design and manage computational pipelines.

\hspace{\parindent}\textbf{Progress to Date:} The project has successfully completed the workflow design phase, including visual workflow design capability through React Flow integration, support for multiple programming languages (Python, JavaScript, Go, Ruby), persistent storage of workflow definitions, and responsive user interface. The directed acyclic graph representation ensures mathematically sound workflow topology while automatic layout algorithms optimize visual presentation. State management architecture maintains consistency between user interface, API backend, and database persistence.

\hspace{\parindent}\textbf{Technical Progress:} The implemented components demonstrate practical integration of visual programming and polyglot computing in an accessible web platform. The architecture balances simplicity with extensibility, positioning the project well for the execution orchestration phase. Type-safe development using TypeScript and Python type hints has proven valuable for code quality and maintainability. The modular design allows the execution engine to be integrated without major architectural changes.

\hspace{\parindent}\textbf{Current Limitations:} As an interim deliverable, the current implementation focuses on workflow design and persistence. The execution orchestration component is not yet implemented. Known limitations include:

\begin{itemize}
    \item No actual container execution - workflows can be designed but not executed
    \item No user authentication or multi-user support
    \item Limited error handling for edge cases
    \item No workflow execution history or logging
    \item Performance testing incomplete for very large workflows (100+ nodes)
\end{itemize}

\hspace{\parindent}These limitations are expected for an interim deliverable and will be addressed in the final implementation phase.

\hspace{\parindent}\textbf{Expected Impact:} Upon completion, the project aims to demonstrate feasibility of creating accessible workflow tools without sacrificing power or flexibility. The interim progress validates the core design decisions and technical approach. The remaining implementation phase will determine whether the platform successfully reduces technical barriers to workflow orchestration for scientific researchers.

\hspace{\parindent}\textbf{Lessons Learned:} The development experience reinforced importance of iterative development, starting with minimal viable product and incrementally adding features. Early user feedback (through supervisor consultation) proved invaluable for interface design decisions. Balancing ideal architecture against practical constraints required engineering judgment developed through project experience.

Technical challenges overcome—particularly state synchronization complexity and asynchronous programming paradigms—built confidence and competence applicable to future software projects. The integration of diverse technologies demonstrated breadth of modern full-stack development while revealing areas for continued learning.

\hspace{\parindent}\textbf{Interim Reflections:} The work completed to date represents significant progress in integrating computer science coursework with practical software system development. The workflow design components provide a solid foundation for the execution engine implementation. The project has successfully applied modern web technologies and software engineering principles to a real scientific computing problem.

\hspace{\parindent}\textbf{Confidence in Completion:} Based on the progress to date and the clear architecture established, we are confident the remaining execution orchestration phase can be completed within the project timeline. The modular design allows the execution engine to be developed and tested independently of the workflow design components. The final deliverable will demonstrate a complete, functional workflow orchestration platform for polyglot scientific computing.

\hspace{\parindent}This interim report documents a project in progress. The final report will include complete implementation details, comprehensive validation results, and deployment documentation.

% References
\newpage
\begin{thebibliography}{99}
\addcontentsline{toc}{section}{REFERENCES}

\bibitem{cohen2020scientific}
Cohen, J., Katz, D. S., Barker, M., Hong, N. C., Haines, R., \& Jay, C. (2020). The Four Pillars of Research Software Engineering. \textit{IEEE Software}, 38(1), 97-105.

\bibitem{deelman2015future}
Deelman, E., Vahi, K., Juve, G., Rynge, M., Callaghan, S., Maechling, P. J., ... \& Wenger, K. (2015). Pegasus, a workflow management system for science automation. \textit{Future Generation Computer Systems}, 46, 17-35.

\bibitem{amstutz2016common}
Amstutz, P., Crusoe, M. R., Tijanić, N., Chapman, B., Chilton, J., Heuer, M., ... \& Stojanovic, N. (2016). Common Workflow Language, v1.0. \textit{Specification, Common Workflow Language working group}.

\bibitem{perkel2019workflow}
Perkel, J. M. (2019). Workflow systems turn raw data into scientific knowledge. \textit{Nature}, 573(7772), 149-150.

\bibitem{stodden2016enhancing}
Stodden, V., \& Miguez, S. (2016). Best Practices for Computational Science: Software Infrastructure and Environments for Reproducible and Extensible Research. \textit{Journal of Open Research Software}, 2(1), e21.

\bibitem{liu2015survey}
Liu, J., Pacitti, E., Valduriez, P., \& Mattoso, M. (2015). A Survey of Data-Intensive Scientific Workflow Management. \textit{Journal of Grid Computing}, 13(4), 457-493.

\bibitem{boettiger2015introduction}
Boettiger, C. (2015). An introduction to Docker for reproducible research. \textit{ACM SIGOPS Operating Systems Review}, 49(1), 71-79.

\bibitem{gruening2018practical}
Grüning, B., Chilton, J., Köster, J., Dale, R., Soranzo, N., van den Beek, M., ... \& Blankenberg, D. (2018). Practical Computational Reproducibility in the Life Sciences. \textit{Cell Systems}, 6(6), 631-635.

\bibitem{crusoe2022methods}
Crusoe, M. R., Abeln, S., Iosup, A., Amstutz, P., Chilton, J., Tijanić, N., ... \& Soiland-Reyes, S. (2022). Methods Included: Standardizing Computational Reuse and Portability with the Common Workflow Language. \textit{Communications of the ACM}, 65(6), 54-63.

\bibitem{burnett2016scaling}
Burnett, M., \& Baker, A. (2016). Scaling Up Visual Programming Languages. \textit{IEEE Computer}, 49(3), 45-52.

\bibitem{green1996usability}
Green, T. R., \& Petre, M. (1996). Usability Analysis of Visual Programming Environments: A 'Cognitive Dimensions' Framework. \textit{Journal of Visual Languages \& Computing}, 7(2), 131-174.

\bibitem{meyerovich2013empirical}
Meyerovich, L. A., \& Rabkin, A. S. (2013). Empirical Analysis of Programming Language Adoption. \textit{ACM SIGPLAN Notices}, 48(10), 1-18.

\bibitem{tarjan1972depth}
Tarjan, R. (1972). Depth-First Search and Linear Graph Algorithms. \textit{SIAM Journal on Computing}, 1(2), 146-160.

\bibitem{kahn1962topological}
Kahn, A. B. (1962). Topological Sorting of Large Networks. \textit{Communications of the ACM}, 5(11), 558-562.

\bibitem{coffman1972optimal}
Coffman Jr, E. G., \& Graham, R. L. (1972). Optimal Scheduling for Two-Processor Systems. \textit{Acta Informatica}, 1(3), 200-213.

\bibitem{russell2018workflow}
Russell, M., Taufer, M., Malony, A., \& Jackson, K. R. (2018). Design Patterns for Scientific Workflow Systems. \textit{Proceedings of the 2018 IEEE 14th International Conference on e-Science}, 376-386.

\bibitem{flux2014facebook}
Facebook Open Source. (2014). Flux Application Architecture. Retrieved from https://facebook.github.io/flux/

\bibitem{redux2015abramov}
Abramov, D. (2015). Redux: Predictable State Container for JavaScript Apps. Retrieved from https://redux.js.org/

\bibitem{fielding2000architectural}
Fielding, R. T. (2000). \textit{Architectural Styles and the Design of Network-based Software Architectures} (Doctoral dissertation). University of California, Irvine.

\bibitem{richardson2008restful}
Richardson, L., \& Ruby, S. (2008). \textit{RESTful Web Services}. O'Reilly Media, Inc.

\bibitem{openapi2021specification}
OpenAPI Initiative. (2021). OpenAPI Specification v3.1.0. Retrieved from https://spec.openapis.org/oas/v3.1.0

\bibitem{date2003database}
Date, C. J. (2003). \textit{An Introduction to Database Systems} (8th ed.). Addison-Wesley.

\bibitem{hellerstein2005readings}
Hellerstein, J. M., \& Stonebraker, M. (Eds.). (2005). \textit{Readings in Database Systems} (4th ed.). MIT Press.

\end{thebibliography}

\end{document}